\documentclass{article}
\usepackage[utf8]{inputenc}
\usepackage[T1]{fontenc}
\usepackage{geometry}
\geometry{
    a4paper,
    left=2.5cm,
    right=2.5cm,
    top=2.5cm,
    bottom=2.5cm
}
\usepackage{graphicx}
\usepackage{float}
\usepackage{hyperref}
\usepackage[english]{babel}
\usepackage{amsmath}
\usepackage{amssymb}
\usepackage{array}
\usepackage{booktabs}
\usepackage{caption}
\usepackage{subcaption}
\usepackage{enumitem}
\usepackage{fancyhdr}
\usepackage{titlesec}
\usepackage{url}
\usepackage{textcomp}
\usepackage{xcolor}
\usepackage{listings}
\usepackage{multirow}
\usepackage{longtable}
\usepackage{tabularx}
\usepackage{times}

% Set up page style
\pagestyle{fancy}
\fancyhf{}
\rfoot{\thepage}
\renewcommand{\headrulewidth}{0pt}

% Title formatting
\titleformat{\chapter}[display]{\normalfont\bfseries}{\chaptertitle\thechapter}{0pt}{\Huge}
\titlespacing*{\chapter}{0pt}{50pt}{40pt}

% Document info
\title{Template Document}
\author{Academic Support Agent}
\date{\today}

\begin{document}
\maketitle

% Para for ar que elementos pr -textuais (da capa at  o sum rio) sejam impressos no anverso da folha
\setboolean{@twoside}{false}

\autor{Nome do aluno}

% Coloque o t tulo em caixa alta.   o padr o da PUC.
% V  no arquivo abnt-PPGInf-PUCMG.sty e procure por esse t tulo (linha 575). Altere para o seu t tulo em caixa alta. Isso ser  utilizado na folha de aprova o.
\titulo{T TULO DA DISSERTA O}

\orientador[Orientador:]{Prof. Dr. Nome}

% Se n o tiver, co-orientador, comente a pr xima linha.
%\coorientador[Co-orientador:]{Professor}

% Texto
\comentario{Disserta o apresentada ao Programa de P s-Gradua o em Inform tica da Pontif cia Universidade Cat lica de Minas Gerais,
como requisito parcial para obten o do t tulo de Mestre em Inform tica.}


% Institui o
\instituicao{Mestrado em Inform tica \par Instituto de Inform tica \par Pontif cia Universidade Cat lica de Minas Gerais}

% Local
\local{Belo Horizonte}

% Data
\data{}
\capa
%Para for ar que a ficha catalogr fica seja impressa no verso da folha de aprova o
\setboolean{@twoside}{true}

% Gera a folha de rosto
\folhaderosto

% Ficha catalogr fica
% Content from pre-texto/ficha-catalografica would be included here


% Para for ar que elementos pr -textuais (da capa at  o sum rio) sejam impressos no anverso da folha
\setboolean{@twoside}{false}

% Folha de aprova o
% Content from pre-texto/folha-aprovacao would be included here


% Dedicat ria
% Content from pre-texto/dedicatoria would be included here


% Agradecimentos
% Content from pre-texto/agradecimentos would be included here


% Ep grafe
% Content from pre-texto/epigrafe would be included here


% Resumo
% Content from pre-texto/resumo would be included here


% Abstract
% Content from pre-texto/abstract would be included here


\makeatletter
\renewcommand\numberline[1]{
	\leftskip 0em
	\rightskip 1.6em
	\parfillskip -\rightskip
	\parindent 0em
	\@tempdima 2.0em
	\vspace{0em} \advance\leftskip \@tempdima \null\nobreak\hskip -\leftskip
	FIGURA \normalfont #1 -- }
\makeatother

% Lista de figuras


\makeatletter
\renewcommand\numberline[1]{
	\leftskip 0em
	\rightskip 1.6em
	\parfillskip -\rightskip
	\parindent 0em
	\@tempdima 2.0em
	\vspace{0em} \advance\leftskip \@tempdima \null\nobreak\hskip -\leftskip
	TABELA \normalfont #1 -- }
\makeatother

% Lista de tabelas


\makeatletter
\renewcommand\numberline[1]{
	\leftskip 0em
	\rightskip 1.6em
	\parfillskip -\rightskip
	\parindent 0em
	\@tempdima 2.0em
	\vspace{0.5em} \advance\leftskip \@tempdima \null\nobreak\hskip -\leftskip
	GR FICO \normalfont #1 -- }
\makeatother

% Lista de graficos
\listof{grafico}{Lista de Graficos}

% Lista de siglas
% Content from pre-texto/lista-siglas.tex would be included here


\makeatletter
\renewcommand\numberline[1]{#1\hspace{0.8em}}
\makeatother

% Altera para espa amento simples a partir daqui
\singlespacing

% Sum rio


% Altera para espa amento 1,5 a partir daqui
\onehalfspacing

%% TEXTUAIS 
% Para for ar que elementos textuais e p s-textuais sejam impressos no anverso e verso das folhas
\setboolean{@twoside}{true}
% Altere o n mero da p gina para o correto. Conte todas as p ginas frente e verso, menos a capa, nclusive a ficha catalogr fica at  a p gina do primeiro cap tulo.
\setcounter{page}{25}

% Cap tulos
% Para for ar que o cap tulo de introdu o comece no anverso
\setboolean{@openright}{true}
% Content from texto/cap1 would be included here

% Os demais cap tulos n o precisam come ar no anverso
\setboolean{@openright}{false}
% Content from texto/cap2 would be included here

% Content from texto/cap3 would be included here


% P S-TEXTUAIS %%
% Bibliografia no arquivo 'Dissertacao.bib'
% Alterar o t tulo das refer ncias para somente 'Refer ncias'
\renewcommand{\bibname}{Refer ncias}
%
%

% Para for ar que os ap ndices e anexos comecem no anverso
%\setboolean{@openright}{true}

\apendice
\begin{apendice}
% Content from pos-texto/apendiceA would be included here

\end{apendice}

\anexo
% Content from pos-texto/anexo would be included here
\end{document}